\documentclass[11pt,a4paper]{article}
\usepackage[utf8]{inputenc}
\usepackage[T1]{fontenc}
\usepackage[catalan]{babel}
\usepackage{lmodern}
\usepackage[sfdefault,lining]{FiraSans}
\usepackage[a4paper,top=2.3cm,bottom=2.2cm,left=1.8cm,right=1.8cm,headheight=1.3cm,headsep=0.4cm,footskip=1cm]{geometry}
\usepackage{amsmath,amssymb}
\usepackage{icomma}
\usepackage{graphicx}
\usepackage[table]{xcolor}
\usepackage{tikz}
\usepackage{fontawesome5}
\usepackage{ccicons}
\usepackage{fancyhdr}
\usepackage{lastpage}
\usepackage{tasks}
\usepackage{needspace}
\usepackage[most]{tcolorbox}
\usepackage{array}
\usepackage[hidelinks,pdftitle={Solucions de la fitxa de la unitat 5: Proporcionalitat i percentatges. Nombres decimals. 1r d'ESO},pdfsubject={Matemàtiques 1r ESO},pdfkeywords={nombres decimals, aproximació, operacions amb decimals, decimals i fraccions, raons, proporcions, proporcionalitat directa, proporcionalitat inversa, regla de tres, percentatges}]{hyperref}

\definecolor{navy}{HTML}{1B2748}
\definecolor{accent}{HTML}{E8603A}
\definecolor{teal}{HTML}{1F9E90}
\definecolor{soft}{HTML}{EEF1F8}
\definecolor{muted}{HTML}{5D6685}

\setlength{\parindent}{0pt}
\setlength{\parskip}{3pt}
\settasks{label=\alph*), label-format=\bfseries\color{navy}, label-width=1.3em, item-indent=1.8em, column-sep=0.6em, after-item-skip=5pt, before-skip=3pt, after-skip=4pt}

% ---------- capçalera i peu ----------
\newcommand{\llicencia}{\href{https://creativecommons.org/licenses/by-nc-sa/4.0/deed.ca}{\ccbyncsa}\ \ Llicència Creative Commons Reconeixement-NoComercial-CompartirIgual 4.0 Internacional (CC BY-NC-SA 4.0)}
\pagestyle{fancy}
\fancyhf{}
\fancyhead[L]{\raisebox{-0.35\height}{\includegraphics[height=0.95cm]{logo-gva.png}}}
\fancyhead[C]{\small\color{muted}Solucions de la fitxa · Unitat 5 · 1r d'ESO}
\fancyhead[R]{\raisebox{-0.35\height}{\includegraphics[height=0.95cm]{logo-vermellar.png}}}
\fancyfoot[L]{\scriptsize\color{muted}\llicencia}
\fancyfoot[R]{\small\color{muted}\thepage}
\renewcommand{\headrulewidth}{0.4pt}
\renewcommand{\headrule}{\color{muted!40}\hrule height \headrulewidth}
\fancypagestyle{primera}{\fancyhf{}\renewcommand{\headrulewidth}{0pt}
  \fancyfoot[L]{\scriptsize\color{muted}\llicencia}
  \fancyfoot[R]{\small\color{muted}\thepage}}

% ---------- nivells i exercicis ----------
\newcommand{\estrelles}[1]{{\color{accent}\ifcase#1\or\faStar\faStar[regular]\faStar[regular]\or\faStar\faStar\faStar[regular]\or\faStar\faStar\faStar\fi}}
\newcounter{exer}
\newcommand{\ex}[1]{\par\needspace{4\baselineskip}\medskip\stepcounter{exer}\noindent{\color{navy}\bfseries\theexer.}\ \ {\small\estrelles{#1}}\ \ }

\newcommand{\seccio}[2]{\par\needspace{10\baselineskip}\bigskip
  \begin{tcolorbox}[colback=navy,colframe=navy,coltext=white,arc=2mm,boxrule=0pt,left=3mm,right=3mm,top=1.2mm,bottom=1.2mm]
  {\large\bfseries #1\quad #2}
  \end{tcolorbox}\vspace{-2pt}}

\newcolumntype{C}[1]{>{\centering\arraybackslash}p{#1}}

% quadre de recordatori (només a la fitxa d'exercicis)
\newcommand{\recorda}[1]{\par\smallskip\begin{tcolorbox}[colback=accent!7,colframe=accent!55,arc=2mm,boxrule=0.5pt,left=2.5mm,right=2.5mm,top=1mm,bottom=1mm]\small{\color{accent}\faLightbulb[regular]}\ \textbf{Recorda:}\ #1\end{tcolorbox}}
\newcommand{\res}[1]{\boxed{#1}}\newcommand{\rest}[1]{\fbox{#1}}

% =====================================================================
% Capçalera FentMates de les fitxes: un color per a cada curs.
% Els colors són els de la pàgina inicial (LEVELS a assets/js/materials.js);
% el logo ve de _eines/logo/pdf/ (python3 _eines/logo/genera.py).
%
% Ús, al final del preàmbul de la unitat:
%   % =====================================================================
% Capçalera FentMates de les fitxes: un color per a cada curs.
% Els colors són els de la pàgina inicial (LEVELS a assets/js/materials.js);
% el logo ve de _eines/logo/pdf/ (python3 _eines/logo/genera.py).
%
% Ús, al final del preàmbul de la unitat:
%   % =====================================================================
% Capçalera FentMates de les fitxes: un color per a cada curs.
% Els colors són els de la pàgina inicial (LEVELS a assets/js/materials.js);
% el logo ve de _eines/logo/pdf/ (python3 _eines/logo/genera.py).
%
% Ús, al final del preàmbul de la unitat:
%   \input{../comu/fentmates.tex}
%   \fmcurs{1eso}
% i, al començament del document:
%   \capcalera{Matemàtiques}{1r ESO}{Fitxa d'exercicis · Unitat 2}{Divisibilitat}
% Opcional abans de \input: \def\fmneed{14} (espai mínim abans de cada secció).
% =====================================================================
\makeatletter
\def\fm@logodir{../../logo/pdf}
\providecommand{\fmneed}{10}
\def\fm@cursnom{}

% clau del curs, color i color del text damunt del color
\newcommand{\fm@curs}[3]{\@namedef{fm@color@#1}{#2}\@namedef{fm@text@#1}{#3}}
\fm@curs{1eso}{E6B84A}{navy}
\fm@curs{2eso}{3B82F6}{white}
\fm@curs{3eso}{FF7A50}{white}
\fm@curs{4eso-a}{A78BFA}{white}
\fm@curs{4eso-b}{6D28D9}{white}
\fm@curs{1batx-mi}{9B1C1C}{white}
\fm@curs{1batx-ccss}{EF6B6B}{white}
\fm@curs{2batx-mii}{1E6B47}{white}
\fm@curs{2batx-ccss}{2BB3A3}{white}

\newcommand{\fmcurs}[1]{%
  \@ifundefined{fm@color@#1}{\PackageError{fentmates}{Curs desconegut: #1}{}}{}%
  \edef\fm@tmp{\noexpand\definecolor{curs}{HTML}{\@nameuse{fm@color@#1}}%
               \noexpand\colorlet{curstext}{\@nameuse{fm@text@#1}}}\fm@tmp
  \colorlet{cursfosc}{curs!72!black}%
  \def\fm@logo{\fm@logodir/vertical-#1.pdf}%
  \def\fm@marca{\fm@logodir/marca-#1.pdf}}

% ---------- capçalera de la primera pàgina ----------
% #1 matèria (petit), #2 curs, #3 tipus de document i unitat, #4 títol
\newcommand{\capcalera}[4]{%
  \gdef\fm@cursnom{#2}%
  \par\noindent
  \begin{tcolorbox}[colback=curs,colframe=curs,boxrule=0pt,arc=2.2mm,width=3.5cm,height=2.15cm,
      valign=center,left=3mm,right=2mm,top=1mm,bottom=1mm,nobeforeafter,box align=center]
    \color{curstext}{\footnotesize #1}\par\vspace{1pt}{\fontsize{19}{21}\selectfont\bfseries #2\par}%
  \end{tcolorbox}\hspace{5mm}%
  \begin{minipage}[c]{\dimexpr\linewidth-3.5cm-5mm-2.9cm\relax}
    {\color{cursfosc}\small\bfseries\MakeUppercase{#3}}\par\vspace{2pt}
    {\color{navy}\fontsize{20}{23}\selectfont\bfseries #4\par}
  \end{minipage}\hfill
  \begin{minipage}[c]{2.6cm}\raggedleft\includegraphics[width=2.6cm]{\fm@logo}\end{minipage}\par
  \vspace{2pt}}

% ---------- barres de secció del color del curs ----------
\renewcommand{\seccio}[2]{\par\needspace{\fmneed\baselineskip}\bigskip
  \begin{tcolorbox}[colback=curs,colframe=curs,coltext=curstext,arc=2mm,boxrule=0pt,left=3mm,right=3mm,top=1.2mm,bottom=1.2mm]
  {\large\bfseries #1\quad #2}
  \end{tcolorbox}\vspace{-2pt}}

% ---------- estrelles de dificultat del color del curs (a Batxillerat no se'n posen) ----------
\@ifundefined{estrelles}{}{%
  \renewcommand{\estrelles}[1]{{\color{curs}\ifcase#1\or\faStar\faStar[regular]\faStar[regular]\or\faStar\faStar\faStar[regular]\or\faStar\faStar\faStar\fi}}}

% ---------- capçalera i peu de les pàgines ----------
% primera pàgina: logos institucionals al peu
\fancypagestyle{primera}{\fancyhf{}\renewcommand{\headrulewidth}{0pt}
  \fancyfoot[L]{\raisebox{-0.32\height}{\includegraphics[height=0.78cm]{logo-gva.png}}\hspace{5mm}%
    \raisebox{-0.32\height}{\includegraphics[height=0.78cm]{logo-vermellar.png}}}
  \fancyfoot[R]{\scriptsize\color{muted}\href{https://creativecommons.org/licenses/by-nc-sa/4.0/deed.ca}{\ccbyncsa}\ \ CC BY-NC-SA 4.0 · fentmates.org\qquad\small\thepage}}
% la resta: el símbol a l'esquerra i el curs a la dreta, amb el títol al centre com abans
\fancyhead[L]{\raisebox{-0.3\height}{\includegraphics[height=0.62cm]{\fm@marca}}}
\fancyhead[R]{\small\bfseries\color{cursfosc}\fm@cursnom}
\renewcommand{\headrule}{\color{curs!45}\hrule height \headrulewidth}
\makeatother

%   \fmcurs{1eso}
% i, al començament del document:
%   \capcalera{Matemàtiques}{1r ESO}{Fitxa d'exercicis · Unitat 2}{Divisibilitat}
% Opcional abans de \input: \def\fmneed{14} (espai mínim abans de cada secció).
% =====================================================================
\makeatletter
\def\fm@logodir{../../logo/pdf}
\providecommand{\fmneed}{10}
\def\fm@cursnom{}

% clau del curs, color i color del text damunt del color
\newcommand{\fm@curs}[3]{\@namedef{fm@color@#1}{#2}\@namedef{fm@text@#1}{#3}}
\fm@curs{1eso}{E6B84A}{navy}
\fm@curs{2eso}{3B82F6}{white}
\fm@curs{3eso}{FF7A50}{white}
\fm@curs{4eso-a}{A78BFA}{white}
\fm@curs{4eso-b}{6D28D9}{white}
\fm@curs{1batx-mi}{9B1C1C}{white}
\fm@curs{1batx-ccss}{EF6B6B}{white}
\fm@curs{2batx-mii}{1E6B47}{white}
\fm@curs{2batx-ccss}{2BB3A3}{white}

\newcommand{\fmcurs}[1]{%
  \@ifundefined{fm@color@#1}{\PackageError{fentmates}{Curs desconegut: #1}{}}{}%
  \edef\fm@tmp{\noexpand\definecolor{curs}{HTML}{\@nameuse{fm@color@#1}}%
               \noexpand\colorlet{curstext}{\@nameuse{fm@text@#1}}}\fm@tmp
  \colorlet{cursfosc}{curs!72!black}%
  \def\fm@logo{\fm@logodir/vertical-#1.pdf}%
  \def\fm@marca{\fm@logodir/marca-#1.pdf}}

% ---------- capçalera de la primera pàgina ----------
% #1 matèria (petit), #2 curs, #3 tipus de document i unitat, #4 títol
\newcommand{\capcalera}[4]{%
  \gdef\fm@cursnom{#2}%
  \par\noindent
  \begin{tcolorbox}[colback=curs,colframe=curs,boxrule=0pt,arc=2.2mm,width=3.5cm,height=2.15cm,
      valign=center,left=3mm,right=2mm,top=1mm,bottom=1mm,nobeforeafter,box align=center]
    \color{curstext}{\footnotesize #1}\par\vspace{1pt}{\fontsize{19}{21}\selectfont\bfseries #2\par}%
  \end{tcolorbox}\hspace{5mm}%
  \begin{minipage}[c]{\dimexpr\linewidth-3.5cm-5mm-2.9cm\relax}
    {\color{cursfosc}\small\bfseries\MakeUppercase{#3}}\par\vspace{2pt}
    {\color{navy}\fontsize{20}{23}\selectfont\bfseries #4\par}
  \end{minipage}\hfill
  \begin{minipage}[c]{2.6cm}\raggedleft\includegraphics[width=2.6cm]{\fm@logo}\end{minipage}\par
  \vspace{2pt}}

% ---------- barres de secció del color del curs ----------
\renewcommand{\seccio}[2]{\par\needspace{\fmneed\baselineskip}\bigskip
  \begin{tcolorbox}[colback=curs,colframe=curs,coltext=curstext,arc=2mm,boxrule=0pt,left=3mm,right=3mm,top=1.2mm,bottom=1.2mm]
  {\large\bfseries #1\quad #2}
  \end{tcolorbox}\vspace{-2pt}}

% ---------- estrelles de dificultat del color del curs (a Batxillerat no se'n posen) ----------
\@ifundefined{estrelles}{}{%
  \renewcommand{\estrelles}[1]{{\color{curs}\ifcase#1\or\faStar\faStar[regular]\faStar[regular]\or\faStar\faStar\faStar[regular]\or\faStar\faStar\faStar\fi}}}

% ---------- capçalera i peu de les pàgines ----------
% primera pàgina: logos institucionals al peu
\fancypagestyle{primera}{\fancyhf{}\renewcommand{\headrulewidth}{0pt}
  \fancyfoot[L]{\raisebox{-0.32\height}{\includegraphics[height=0.78cm]{logo-gva.png}}\hspace{5mm}%
    \raisebox{-0.32\height}{\includegraphics[height=0.78cm]{logo-vermellar.png}}}
  \fancyfoot[R]{\scriptsize\color{muted}\href{https://creativecommons.org/licenses/by-nc-sa/4.0/deed.ca}{\ccbyncsa}\ \ CC BY-NC-SA 4.0 · fentmates.org\qquad\small\thepage}}
% la resta: el símbol a l'esquerra i el curs a la dreta, amb el títol al centre com abans
\fancyhead[L]{\raisebox{-0.3\height}{\includegraphics[height=0.62cm]{\fm@marca}}}
\fancyhead[R]{\small\bfseries\color{cursfosc}\fm@cursnom}
\renewcommand{\headrule}{\color{curs!45}\hrule height \headrulewidth}
\makeatother

%   \fmcurs{1eso}
% i, al començament del document:
%   \capcalera{Matemàtiques}{1r ESO}{Fitxa d'exercicis · Unitat 2}{Divisibilitat}
% Opcional abans de \input: \def\fmneed{14} (espai mínim abans de cada secció).
% =====================================================================
\makeatletter
\def\fm@logodir{../../logo/pdf}
\providecommand{\fmneed}{10}
\def\fm@cursnom{}

% clau del curs, color i color del text damunt del color
\newcommand{\fm@curs}[3]{\@namedef{fm@color@#1}{#2}\@namedef{fm@text@#1}{#3}}
\fm@curs{1eso}{E6B84A}{navy}
\fm@curs{2eso}{3B82F6}{white}
\fm@curs{3eso}{FF7A50}{white}
\fm@curs{4eso-a}{A78BFA}{white}
\fm@curs{4eso-b}{6D28D9}{white}
\fm@curs{1batx-mi}{9B1C1C}{white}
\fm@curs{1batx-ccss}{EF6B6B}{white}
\fm@curs{2batx-mii}{1E6B47}{white}
\fm@curs{2batx-ccss}{2BB3A3}{white}

\newcommand{\fmcurs}[1]{%
  \@ifundefined{fm@color@#1}{\PackageError{fentmates}{Curs desconegut: #1}{}}{}%
  \edef\fm@tmp{\noexpand\definecolor{curs}{HTML}{\@nameuse{fm@color@#1}}%
               \noexpand\colorlet{curstext}{\@nameuse{fm@text@#1}}}\fm@tmp
  \colorlet{cursfosc}{curs!72!black}%
  \def\fm@logo{\fm@logodir/vertical-#1.pdf}%
  \def\fm@marca{\fm@logodir/marca-#1.pdf}}

% ---------- capçalera de la primera pàgina ----------
% #1 matèria (petit), #2 curs, #3 tipus de document i unitat, #4 títol
\newcommand{\capcalera}[4]{%
  \gdef\fm@cursnom{#2}%
  \par\noindent
  \begin{tcolorbox}[colback=curs,colframe=curs,boxrule=0pt,arc=2.2mm,width=3.5cm,height=2.15cm,
      valign=center,left=3mm,right=2mm,top=1mm,bottom=1mm,nobeforeafter,box align=center]
    \color{curstext}{\footnotesize #1}\par\vspace{1pt}{\fontsize{19}{21}\selectfont\bfseries #2\par}%
  \end{tcolorbox}\hspace{5mm}%
  \begin{minipage}[c]{\dimexpr\linewidth-3.5cm-5mm-2.9cm\relax}
    {\color{cursfosc}\small\bfseries\MakeUppercase{#3}}\par\vspace{2pt}
    {\color{navy}\fontsize{20}{23}\selectfont\bfseries #4\par}
  \end{minipage}\hfill
  \begin{minipage}[c]{2.6cm}\raggedleft\includegraphics[width=2.6cm]{\fm@logo}\end{minipage}\par
  \vspace{2pt}}

% ---------- barres de secció del color del curs ----------
\renewcommand{\seccio}[2]{\par\needspace{\fmneed\baselineskip}\bigskip
  \begin{tcolorbox}[colback=curs,colframe=curs,coltext=curstext,arc=2mm,boxrule=0pt,left=3mm,right=3mm,top=1.2mm,bottom=1.2mm]
  {\large\bfseries #1\quad #2}
  \end{tcolorbox}\vspace{-2pt}}

% ---------- estrelles de dificultat del color del curs (a Batxillerat no se'n posen) ----------
\@ifundefined{estrelles}{}{%
  \renewcommand{\estrelles}[1]{{\color{curs}\ifcase#1\or\faStar\faStar[regular]\faStar[regular]\or\faStar\faStar\faStar[regular]\or\faStar\faStar\faStar\fi}}}

% ---------- capçalera i peu de les pàgines ----------
% primera pàgina: logos institucionals al peu
\fancypagestyle{primera}{\fancyhf{}\renewcommand{\headrulewidth}{0pt}
  \fancyfoot[L]{\raisebox{-0.32\height}{\includegraphics[height=0.78cm]{logo-gva.png}}\hspace{5mm}%
    \raisebox{-0.32\height}{\includegraphics[height=0.78cm]{logo-vermellar.png}}}
  \fancyfoot[R]{\scriptsize\color{muted}\href{https://creativecommons.org/licenses/by-nc-sa/4.0/deed.ca}{\ccbyncsa}\ \ CC BY-NC-SA 4.0 · fentmates.org\qquad\small\thepage}}
% la resta: el símbol a l'esquerra i el curs a la dreta, amb el títol al centre com abans
\fancyhead[L]{\raisebox{-0.3\height}{\includegraphics[height=0.62cm]{\fm@marca}}}
\fancyhead[R]{\small\bfseries\color{cursfosc}\fm@cursnom}
\renewcommand{\headrule}{\color{curs!45}\hrule height \headrulewidth}
\makeatother

\fmcurs{1eso}

\begin{document}
\thispagestyle{primera}

\capcalera{Matemàtiques}{1r ESO}{Solucions de la fitxa · Unitat 5}{Proporcionalitat i percentatges. Nombres decimals}
\vspace{8pt}

\begin{tcolorbox}[colback=soft,colframe=soft,arc=2mm,left=3mm,right=3mm,top=1.5mm,bottom=1.5mm]
\small Intenta resoldre cada exercici abans de mirar-ne la solució. Si no t'ix el mateix resultat, repassa els passos o consulta els apunts interactius de la unitat. Recorda que les fraccions s'han de donar en forma irreductible i els preus, arrodonits als cèntims.
\end{tcolorbox}

% =====================================================
\seccio{1}{Nombres decimals}

\ex{1} \begin{tasks}(2)
  \task $7,25$: 7 unitats i 25 centèsimes
  \task $0,408$: 408 mil·lèsimes
  \task $12,06$: 12 unitats i 6 centèsimes
  \task $3,5$: 3 unitats i 5 dècimes
  \task $100,001$: 100 unitats i 1 mil·lèsima
  \task $0,09$: 9 centèsimes
\end{tasks}

\ex{1} \begin{tasks}(3)
  \task $\res{3,045}$
  \task $\res{0,12}$
  \task $\res{20,9}$
  \task $\res{0,007}$
  \task $\res{100,05}$
  \task $\res{4,300}$
\end{tasks}

\ex{1} \begin{tasks}(2)
  \task El 7 en $34,571$: 7 centèsimes = 0,07
  \task El 3 en $2,738$: 3 centèsimes = 0,03
  \task El 7 en $607,15$: 7 unitats
  \task El 8 en $0,0809$: 8 centèsimes = 0,08
  \task El 4 en $15,4$: 4 dècimes = 0,4
  \task El 6 en $9,206$: 6 mil·lèsimes = 0,006
\end{tasks}

\ex{2} \begin{tasks}(3)
  \task $\res{10}$ dècimes
  \task $\res{10}$ centèsimes
  \task $\res{300}$ centèsimes
  \task $\res{0,45}$ unitats
  \task $\res{2,5}$ dècimes
  \task $\res{700}$ mil·lèsimes
\end{tasks}

\ex{2} \begin{tasks}(1)
  \task $45,312=40+5+0,3+0,01+0,002$
  \task $7,09=7+0,09$
  \task $120,506=100+20+0,5+0,006$
\end{tasks}

\ex{3} a) $\res{40,602}$ \quad b) $\res{803,05}$ \quad c) $\res{0,153}$. \ Són iguals a $0,15$: $\res{0,150}$ i $\res{0,1500}$ (els zeros finals no canvien el valor).


% =====================================================
\seccio{2}{Comparar, representar i aproximar decimals}

\ex{1} \begin{tasks}(3)
  \task $4,7\ \res{>}\ 4,68$
  \task $0,09\ \res{<}\ 0,1$
  \task $5,307\ \res{<}\ 5,37$
  \task $12,5\ \res{=}\ 12,50$
  \task $3,099\ \res{<}\ 3,1$
  \task $0,45\ \res{>}\ 0,405$
\end{tasks}

\ex{1} $\res{2,0<2,05<2,15<2,5<2,505<2,55}$

\ex{2} A $=\res{3,2}$, \ B $=\res{3,7}$, \ C $=\res{3,9}$, \ D $=\res{0,53}$, \ E $=\res{0,58}$.

\ex{2} Dividim la unitat en 20 parts (cada una val $0,05$):\par\smallskip{\centering \begin{tikzpicture}[baseline=0]\draw[-latex,thick] (-0.3,0)--(12.50,0);\draw[thick] (0.000,0.20)--(0.000,-0.20);\node[below=2pt,font=\small] at (0.000,-0.1) {$1$};\draw[muted] (0.600,0.12)--(0.600,-0.12);\draw[muted] (1.200,0.12)--(1.200,-0.12);\draw[muted] (1.800,0.12)--(1.800,-0.12);\draw[muted] (2.400,0.12)--(2.400,-0.12);\draw[muted] (3.000,0.12)--(3.000,-0.12);\draw[muted] (3.600,0.12)--(3.600,-0.12);\draw[muted] (4.200,0.12)--(4.200,-0.12);\draw[muted] (4.800,0.12)--(4.800,-0.12);\draw[muted] (5.400,0.12)--(5.400,-0.12);\draw[muted] (6.000,0.12)--(6.000,-0.12);\draw[muted] (6.600,0.12)--(6.600,-0.12);\draw[muted] (7.200,0.12)--(7.200,-0.12);\draw[muted] (7.800,0.12)--(7.800,-0.12);\draw[muted] (8.400,0.12)--(8.400,-0.12);\draw[muted] (9.000,0.12)--(9.000,-0.12);\draw[muted] (9.600,0.12)--(9.600,-0.12);\draw[muted] (10.200,0.12)--(10.200,-0.12);\draw[muted] (10.800,0.12)--(10.800,-0.12);\draw[muted] (11.400,0.12)--(11.400,-0.12);\draw[thick] (12.000,0.20)--(12.000,-0.20);\node[below=2pt,font=\small] at (12.000,-0.1) {$2$};\fill[accent] (4.800,0) circle (0.09) node[above=3pt,text=accent,font=\small]{$1,4$};\fill[accent] (9.000,0) circle (0.09) node[above=3pt,text=accent,font=\small]{$1,75$};\end{tikzpicture}\par}

\ex{2} \begin{tasks}(3)
  \task Per exemple $0,31$; $0,35$; $0,39$
  \task Per exemple $2,711$; $2,715$; $2,719$
  \task Per exemple $5,01$; $5,05$; $5,09$
\end{tasks}\par (Hi ha infinites respostes.)

\ex{2} \par\smallskip{\centering\renewcommand{\arraystretch}{1.35}\begin{tabular}{|c|C{2.3cm}|C{2.3cm}|C{2.3cm}|}\hline Nombre & Truncat a les centèsimes & Arrodonit a les dècimes & Arrodonit a les centèsimes\\\hline $3,4682$ & $3,46$ & $\res{3,5}$ & $\res{3,47}$\\\hline $15,8962$ & $15,89$ & $\res{15,9}$ & $\res{15,90}$\\\hline $0,7049$ & $0,70$ & $\res{0,7}$ & $\res{0,70}$\\\hline $29,996$ & $29,99$ & $\res{30,0}$ & $\res{30,00}$\\\hline\end{tabular}\par}

\ex{3} Arrodonim a cèntims: $\res{12,46\text{\,€}}$. \ Per exemple $\res{2,47}$ (qualsevol des de $2,45$ fins a $2,5$ sense arribar-hi). \ El més xicotet és $\res{7,50}$.


% =====================================================
\seccio{3}{Sumar, restar i multiplicar decimals}

\ex{1} \begin{tasks}(3)
  \task $23,45+7,8=\res{31,25}$
  \task $40-12,35=\res{27,65}$
  \task $0,9+0,35=\res{1,25}$
  \task $15-6,47=\res{8,53}$
  \task $3,06+12,9+0,4=\res{16,36}$
  \task $100-0,01=\res{99,99}$
\end{tasks}

\ex{1} \begin{tasks}(3)
  \task $4,6\cdot 3=\res{13,8}$
  \task $2,35\cdot 0,4=\res{0,94}$
  \task $0,6\cdot 0,05=\res{0,03}$
  \task $12,5\cdot 1,2=\res{15}$
  \task $3,25\cdot 1,4=\res{4,55}$
  \task $0,07\cdot 0,3=\res{0,021}$
\end{tasks}

\ex{1} \begin{tasks}(4)
  \task $3,7\cdot 100=\res{370}$
  \task $0,052\cdot 1000=\res{52}$
  \task $45,2:1000=\res{0,0452}$
  \task $8:100=\res{0,08}$
  \task $6,15\cdot 10=\res{61,5}$
  \task $0,3:10=\res{0,03}$
  \task $250\cdot 0,1=\res{25}$
  \task $1,4\cdot 0,01=\res{0,014}$
\end{tasks}

\ex{2} $3,8\cdot2,5=\res{9,50}$ (prop de $4\cdot2=8$); \ $0,45\cdot1,2=\res{0,540}$ (prop de $0,5$); \ $12,6\cdot0,03=\res{0,378}$ (2+1 decimals sumen 3).

\ex{2} \begin{tasks}(3)
  \task $\res{2,25}$
  \task $\res{5,25}$
  \task $\res{100}$
\end{tasks}

\ex{2} Pomes: $2,5\cdot1,80=4,50$\,€; \ cireres: $1,5\cdot6,40=9,60$\,€; \ total $14,10$\,€. \ Ens tornen $20-14,10=\res{5,90\text{\,€}}$.

\ex{3} Laia: $25\cdot8=200$ m; \ germà: $25\cdot5,5=137,5$ m; \ diferència $\res{62,5\text{ m}}$. \ Temps: $8\cdot0,45=\res{3,6\text{ minuts}}$.

\ex{3} Tercer costat: $17,25-(4,8+6,07)=17,25-10,87=\res{6,38\text{ cm}}$. \ Perímetre: $17,25\cdot10=\res{172,5\text{ cm}}$.


% =====================================================
\seccio{4}{Dividir amb decimals}

\ex{1} \begin{tasks}(3)
  \task $8,4:4=\res{2,1}$
  \task $9,75:3=\res{3,25}$
  \task $6,24:3=\res{2,08}$
  \task $0,8:2=\res{0,4}$
  \task $15,3:5=\res{3,06}$
  \task $1,26:6=\res{0,21}$
\end{tasks}

\ex{1} \begin{tasks}(3)
  \task $7:4=\res{1,75}$
  \task $15:8=\res{1,875}$
  \task $9:4=\res{2,25}$
  \task $3:5=\res{0,6}$
  \task $1:8=\res{0,125}$
  \task $21:12=\res{1,75}$
\end{tasks}

\ex{2} \begin{tasks}(3)
  \task $4,5:0,9=45:9=\res{5}$
  \task $2,73:0,7=27,3:7=\res{3,9}$
  \task $6:0,12=600:12=\res{50}$
  \task $12,6:0,35=1260:35=\res{36}$
  \task $1,44:1,2=14,4:12=\res{1,2}$
  \task $0,9:0,03=90:3=\res{30}$
\end{tasks}

\ex{2} a) $\res{3,33}$, residu $0,01$ \quad b) $\res{0,71}$, residu $0,03$ \quad c) $\res{4,23}$, residu $0,02$.

\ex{2} $1,5:0,25=150:25=\res{6\text{ gots}}$; \ $5:0,25=500:25=\res{20\text{ gots}}$.

\ex{3} $11,39:0,85=1139:85=\res{13,40\text{\,€/kg}}$. \ $1,2\cdot13,40=\res{16,08\text{\,€}}$.

\ex{3} $31,5:4,5=\res{7\text{ L}}$ cada 100 km. \ Cost: $31,5\cdot1,56=\res{49,14\text{\,€}}$.


% =====================================================
\seccio{5}{Decimals i fraccions}

\ex{1} \begin{tasks}(4)
  \task $0,6=\dfrac{6}{10}=\res{\dfrac{3}{5}}$
  \task $0,45=\dfrac{45}{100}=\res{\dfrac{9}{20}}$
  \task $1,2=\dfrac{12}{10}=\res{\dfrac{6}{5}}$
  \task $0,004=\dfrac{4}{1000}=\res{\dfrac{1}{250}}$
  \task $0,125=\dfrac{125}{1000}=\res{\dfrac{1}{8}}$
  \task $2,75=\dfrac{275}{100}=\res{\dfrac{11}{4}}$
  \task $0,08=\dfrac{8}{100}=\res{\dfrac{2}{25}}$
  \task $3,5=\dfrac{35}{10}=\res{\dfrac{7}{2}}$
\end{tasks}

\ex{1} \begin{tasks}(4)
  \task $\res{0,7}$
  \task $\res{0,35}$
  \task $\res{0,004}$
  \task $\res{1,25}$
\end{tasks}

\ex{2} \begin{tasks}(2)
  \task $\dfrac{9}{20}=\res{0,45}$ \ (exacte)
  \task $\dfrac{5}{9}=\res{0,\overline{5}}$ \ (periòdic pur)
  \task $\dfrac{7}{12}=\res{0,58\overline{3}}$ \ (periòdic mixt)
  \task $\dfrac{3}{8}=\res{0,375}$ \ (exacte)
  \task $\dfrac{4}{11}=\res{0,\overline{36}}$ \ (periòdic pur)
  \task $\dfrac{1}{6}=\res{0,1\overline{6}}$ \ (periòdic mixt)
  \task $\dfrac{13}{5}=\res{2,6}$ \ (exacte)
  \task $\dfrac{2}{3}=\res{0,\overline{6}}$ \ (periòdic pur)
\end{tasks}

\ex{2} $\dfrac{2}{3}=0,\overline{6}$ \ i \ $\dfrac{3}{4}=0,75$: \qquad $0,6<\res{\dfrac{2}{3}}<0,67<0,7<\dfrac{3}{4}$

\ex{2} \begin{tasks}(3)
  \task $\dfrac{5}{8}=0,625>0,62$
  \task $\dfrac{1}{3}=0,333\ldots>0,33$
  \task $\dfrac{7}{5}=1,4$: són iguals
\end{tasks}

\ex{3} $\res{\dfrac{5}{3}}=1,\overline{6}$; \qquad $\res{\dfrac{3}{8}}=0,375$.

\ex{3} Marta: $0,25=\dfrac{2}{8}$ (2 porcions); Pau: $\dfrac{3}{8}=0,375$. \ \rest{El Pau} n'ha menjat més. \ Queda $1-\dfrac{2}{8}-\dfrac{3}{8}=\res{\dfrac{3}{8}=0,375}$.


% =====================================================
\seccio{6}{Raons i proporcions}

\ex{1} Pomes/peres: $\dfrac{18}{24}=\res{\dfrac{3}{4}}=0,75$: per cada 3 pomes hi ha 4 peres. \ Peres/total: $\dfrac{24}{42}=\res{\dfrac{4}{7}}$: 4 de cada 7 fruites són peres.

\ex{1} \begin{tasks}(2)
  \task $6\cdot12=72=9\cdot8$: \rest{sí}
  \task $4\cdot15=60=10\cdot6$: \rest{sí}
  \task $2,5\cdot9=22,5\neq4\cdot5=20$: \rest{no}
  \task $3\cdot28=84=7\cdot12$: \rest{sí}
\end{tasks}

\ex{2} \begin{tasks}(3)
  \task $x=\dfrac{5\cdot12}{3}=\res{20}$
  \task $x=\dfrac{8\cdot15}{24}=\res{5}$
  \task $x=\dfrac{4\cdot25}{10}=\res{10}$
  \task $x=\dfrac{2,5\cdot12}{4}=\res{7,5}$
  \task $x=\dfrac{0,3\cdot8}{1,2}=\res{2}$
  \task $x=\dfrac{7\cdot6}{21}=\res{2}$
\end{tasks}

\ex{2} $\dfrac{300}{180}=\dfrac{750}{x}\ \Rightarrow\ x=\dfrac{180\cdot750}{300}=\res{450\text{ mL}}$

\ex{3} $\dfrac{x}{12}=\dfrac{7}{2}\Rightarrow x=\res{42\text{ anys}}$. \ Fa 8 anys: mare 34, filla 4; \ $\dfrac{34}{4}=\dfrac{17}{2}\neq\dfrac{7}{2}$: \rest{la raó no es manté} (les edats no són proporcionals).


% =====================================================
\seccio{7}{Proporcionalitat directa}

\ex{1} Constant: $13:2=\res{6,50}$\,€ per entrada.\par\smallskip{\centering {\renewcommand{\arraystretch}{1.35}\begin{tabular}{|l|C{1.25cm}|C{1.25cm}|C{1.25cm}|C{1.25cm}|C{1.25cm}|}\hline Entrades & 1 & 2 & 4 & $\res{7}$ & 10\\\hline Preu (€) & $\res{6,50}$ & 13 & $\res{26}$ & 45,50 & $\res{65}$\\\hline\end{tabular}}\par}

\ex{1} \begin{tasks}(2)
  \task \rest{Sí}
  \task \rest{No}: augmenten junts però no en la mateixa proporció.
  \task \rest{Sí}: perímetre $=4\cdot$ costat.
  \task \rest{Sí}
\end{tasks}

\ex{1} Un bolígraf: $3,20:4=0,80$\,€. \ 9 bolígrafs: $\res{7,20\text{\,€}}$; \ 15 bolígrafs: $\res{12\text{\,€}}$.

\ex{2} $\begin{array}{rcl}100\text{ km}&\to&6,5\text{ L}\\340\text{ km}&\to&x\end{array}\qquad x=\dfrac{6,5\cdot340}{100}=\res{22,1\text{ L}}$

\ex{2} $x=\dfrac{8\cdot450}{120}=\res{30\text{ min}}$. \ En 30 min, $120:8=15$ L/min, \ $15\cdot30=\res{450\text{ L}}$ (justament el dipòsit).

\ex{2} Multipliquem per $9:6=1,5$: \ $\res{375\text{ g}}$ de farina, \ $\res{6}$ ous \ i \ $\res{0,225\text{ L}}$ de llet.

\ex{3} 1 cm $\to 7,5$ km. \ $9,4\cdot7,5=\res{70,5\text{ km}}$; \ $6:7,5=\res{0,8\text{ cm}}$.

\ex{3} \par\smallskip{\centering {\renewcommand{\arraystretch}{1.35}\begin{tabular}{|l|C{1.25cm}|C{1.25cm}|C{1.25cm}|C{1.25cm}|}\hline km & 1 & 2 & 5 & 10\\\hline Preu (€) & 4,20 & 5,40 & 9 & 15\\\hline\end{tabular}}\par}\smallskip \rest{No és proporcional}: $\dfrac{4,20}{1}\neq\dfrac{5,40}{2}$ (el doble de km no costa el doble).


% =====================================================
\seccio{8}{Proporcionalitat inversa}

\ex{1} Constant: $40\cdot6=\res{240}$ (la distància).\par\smallskip{\centering {\renewcommand{\arraystretch}{1.35}\begin{tabular}{|l|C{1.25cm}|C{1.25cm}|C{1.25cm}|C{1.25cm}|C{1.25cm}|}\hline Velocitat (km/h) & 40 & 60 & 80 & $\res{96}$ & 120\\\hline Temps (h) & 6 & $\res{4}$ & $\res{3}$ & 2,5 & $\res{2}$\\\hline\end{tabular}}\par}

\ex{1} \begin{tasks}(2)
  \task \rest{Inversa}
  \task \rest{Directa}
  \task \rest{Inversa}
  \task \rest{Cap de les dues}
\end{tasks}

\ex{1} Una aixeta: $3\cdot10=30$ h. \ 5 aixetes: $30:5=\res{6\text{ h}}$; \ 2 aixetes: $30:2=\res{15\text{ h}}$.

\ex{2} $\begin{array}{rcl}12\text{ pers.}&\to&30\text{ dies}\\20\text{ pers.}&\to&x\end{array}\qquad x=\dfrac{12\cdot30}{20}=\res{18\text{ dies}}$

\ex{2} $x=\dfrac{90\cdot2}{120}=\res{1,5\text{ h}}$ (1 h 30 min). \ Velocitat: $\dfrac{90\cdot2}{3}=\res{60\text{ km/h}}$.

\ex{3} Inversa: $\dfrac{5\cdot12}{4}=\res{15\text{ màquines}}$. \ Amb 3: $\dfrac{5\cdot12}{3}=\res{20\text{ h}}$.

\ex{3} Capacitat: $60\cdot5=300$ L. \ $300:1,5=\res{200\text{ garrafes}}$; \ $300:25=\res{12\text{ L}}$.


% =====================================================
\seccio{9}{Percentatges}

\ex{1} \par\smallskip{\centering\renewcommand{\arraystretch}{1.5}\begin{tabular}{|c|C{2.4cm}|C{2.4cm}|}\hline Percentatge & Fracció irreductible & Decimal\\\hline $40\,\%$ & $\res{\dfrac{2}{5}}$ & $\res{0,4}$\\\hline $5\,\%$ & $\res{\dfrac{1}{20}}$ & $\res{0,05}$\\\hline $75\,\%$ & $\res{\dfrac{3}{4}}$ & $\res{0,75}$\\\hline $12\,\%$ & $\res{\dfrac{3}{25}}$ & $\res{0,12}$\\\hline $120\,\%$ & $\res{\dfrac{6}{5}}$ & $\res{1,2}$\\\hline $2,5\,\%$ & $\res{\dfrac{1}{40}}$ & $\res{0,025}$\\\hline\end{tabular}\par}

\ex{1} \begin{tasks}(3)
  \task $45\cdot0,2=\res{9}$
  \task $120\cdot0,75=\res{90}$
  \task $60\cdot0,15=\res{9}$
  \task $350\cdot0,08=\res{28}$
  \task $250\cdot0,12=\res{30}$
  \task $60\cdot0,075=\res{4,5}$
\end{tasks}

\ex{1} \begin{tasks}(6)
  \task $\res{35\,\%}$
  \task $\res{6\,\%}$
  \task $\res{150\,\%}$
  \task $\res{60\,\%}$
  \task $\res{35\,\%}$
  \task $\res{12,5\,\%}$
\end{tasks}

\ex{2} $\dfrac{91}{650}=0,14=\res{14\,\%}$; \qquad $\dfrac{6}{25}=0,24=\res{24\,\%}$.

\ex{2} Excel·lent: $28\cdot0,25=\res{7}$; \ notable: $28\cdot0,5=\res{14}$; \ la resta: $100-75=\res{25\,\%}$ (7 alumnes).

\ex{3} $135:0,45=13\,500:45=\res{300\text{ pàgines}}$. \ Comprovació: $300\cdot0,45=135$.

\ex{3} $50\cdot0,3=15=30\cdot0,5$ \ i \ $80\cdot0,2=16=20\cdot0,8$: \rest{són iguals}. \ Sempre passa perquè $\dfrac{30\cdot50}{100}=\dfrac{50\cdot30}{100}$ (l'ordre dels factors no altera el producte).


% =====================================================
\seccio{10}{Augments i disminucions percentuals}

\ex{1} \begin{tasks}(3)
  \task augment del 30 \%: $\res{1,3}$
  \task disminució del 8 \%: $\res{0,92}$
  \task augment del 4 \%: $\res{1,04}$
  \task rebaixa del 50 \%: $\res{0,5}$
  \task IVA del 21 \%: $\res{1,21}$
  \task descompte del 15 \%: $\res{0,85}$
\end{tasks}

\ex{1} \begin{tasks}(2)
  \task $80\cdot0,85=\res{68}$\,€
  \task $50\cdot1,21=\res{60,50}$\,€
  \task $650\cdot0,88=\res{572}$\,€
  \task $18\cdot1,04=\res{18,72}$\,€
  \task $2400\cdot1,05=\res{2520}$\,€
  \task $64,90\cdot0,7=\res{45,43}$\,€
\end{tasks}

\ex{2} Descompte: $5$\,€, \ $\dfrac{5}{25}=0,2=\res{20\,\%}$. \ Pujada: $18$\,€, \ $\dfrac{18}{120}=0,15=\res{15\,\%}$.

\ex{2} Amb IVA: $450\cdot1,21=\res{544,50\text{\,€}}$. \ Amb descompte: $544,50\cdot0,9=\res{490,05\text{\,€}}$.

\ex{3} $108,90:1,21=\res{90\text{\,€}}$; \qquad $68:0,85=\res{80\text{\,€}}$.

\ex{3} $200\cdot1,1=220$; \ $220\cdot0,9=\res{198\text{\,€}}$: no torna a 200\,€. \ Amb $-20\,\%$ i $+20\,\%$: $200\cdot0,8\cdot1,2=\res{192\text{\,€}}$ (un 4\,\% menys).

\ex{3} $2400\cdot1,05=2520$; \ $2520\cdot0,95=\res{2394}$ habitants. \ Variació: $\dfrac{2394}{2400}=0,9975$: \ ha disminuït un $\res{0,25\,\%}$.


\end{document}
